\SourceLeaf{037}
\[
\begin{aligned}
aA+bB+cC+\text{etc.}+l&=\lambda\\
a'A+b'B+c'C+\text{etc.}+l'&=\lambda'\\
a''A+b''B+c''C+\text{etc.}+l''&=\lambda''
\end{aligned}
\]
etc.; statuamus praeterea
\[
\lambda\lambda+\lambda'\lambda'+\lambda''\lambda''+\text{etc.}=M
\]
ita ut sit \(M\) valor functionis \(\Omega\) valoribus maxime plausibilibus indeterminatarum respondens, adeoque per ea, quae in art. 20 demonstravimus, valor minimus huius functionis. Hinc erit \(a\lambda+a'\lambda'+a''\lambda''+\text{etc.}\) valor ipsius \(\xi\), valoribus \(x=A,\allowbreak y=B,\allowbreak z=C\) etc. respondens, adeoque \(=0\), i. e. habebimus
\[
\Sigma a\lambda=0
\]
et perinde fiet
\[
\begin{aligned}
&\Sigma b\lambda=0,\quad\Sigma c\lambda=0\quad\text{etc.;}\\
&\text{nec non}\quad\Sigma\alpha\lambda=0,\quad\Sigma\beta\lambda=0,\quad\Sigma\gamma\lambda=0\quad\text{etc.}
\end{aligned}
\]
Denique multiplicando expressiones ipsarum \(\lambda,\lambda',\lambda''\) etc. per \(\lambda,\lambda',\lambda''\) etc. resp., et addendo, obtinemus \(l\lambda+l'\lambda'+l''\lambda''+\text{etc.}=\lambda\lambda+\lambda'\lambda'+\lambda''\lambda''+\text{etc.}\), sive
\[
\Sigma l\lambda=M
\]
\Article{26}
Substituendo in aequatione \(v=ax+by+cz+\text{etc.}+l\), pro \(x,y,z\) etc. expressiones VII. art. 21, prodibit, adhibitis reductionibus ex praecedentibus obviis,
\[
v=\alpha\xi+\beta\eta+\gamma\zeta+\text{etc.}+\lambda
\]
et perinde erit indefinite
\[
\begin{aligned}
v'&=\alpha'\xi+\beta'\eta+\gamma'\zeta+\text{etc.}+\lambda'\\
v''&=\alpha''\xi+\beta''\eta+\gamma''\zeta+\text{etc.}+\lambda''
\end{aligned}
\]
etc. Multiplicando vel has aequationes, vel aequationes I art. 20 resp. per \(\lambda,\lambda',\lambda''\) etc., et addendo, discimus esse indefinite
\[
\lambda v+\lambda'v'+\lambda''v''+\text{etc.}=M
\]
\SourceLeaf{038}
\Article{27}
Functio \(\Omega\) indefinite in pluribus formis exhiberi potest, quas evolvere operae pretium erit. Ac primo quidem quadrando aequationes I art. 20 et addendo, statim fit
\[
\begin{aligned}
\Omega={}&xx\Sigma aa+yy\Sigma bb+zz\Sigma cc+\text{etc.}\\
&+2xy\Sigma ab+2xz\Sigma ac+2yz\Sigma bc+\text{etc.}\\
&+2x\Sigma al+2y\Sigma bl+2z\Sigma cl+\text{etc.}+\Sigma ll
\end{aligned}
\]
quae est forma \textit{prima}.

Multiplicando easdem aequationes resp. per \(v,v',v''\) etc., et addendo, obtinemus:
\[
\Omega=\xi x+\eta y+\zeta z+\text{etc.}+lv+l'v'+l''v''+\text{etc.}
\]
atque hinc, substituendo pro \(v,v',v''\) etc. expressiones in art. praec. traditas,
\[
\Omega=\xi x+\eta y+\zeta z+\text{etc.}-A\xi-B\eta-C\zeta-\text{etc.}+M
\]
sive
\[
\Omega=\xi(x-A)+\eta(y-B)+\zeta(z-C)+\text{etc.}+M
\]
quae est forma \textit{secunda}.

Substituendo in forma secunda pro \(x-A,y-B,z-C\) etc. expressiones VII. art. 21, obtinemus formam \textit{tertiam}:
\[
\begin{aligned}
\Omega={}&[\alpha\alpha]\xi\xi+[\beta\beta]\eta\eta+[\gamma\gamma]\zeta\zeta+\text{etc.}\\
&+2[\alpha\beta]\xi\eta+2[\alpha\gamma]\xi\zeta+2[\beta\gamma]\eta\zeta+\text{etc.}+M
\end{aligned}
\]
His adiungi potest forma \textit{quarta}, ex forma tertia atque formulis art. praec. sponte demanans:
\[
\Omega=(v-\lambda)^2+(v'-\lambda')^2+(v''-\lambda'')^2+\text{etc.}+M,\quad\text{sive}
\]
\[
\Omega=M+\Sigma(v-\lambda)^2
\]
quae forma conditionem minimi directe ob oculos sistit.
\Article{28}
Sint \(e,e',e''\) etc. errores in observationibus, quae dederunt \(V=L,\allowbreak V'=L'\quad V''=L''\) etc., commissi, i. e. sint valores veri functionum \(V,V',V''\) etc. resp. \(L-e,L'-e',L''-e''\) etc. adeoque valores veri ipsarum \(v,v',v''\) etc. resp. \(-e\BareRoot{p},\allowbreak-e'\BareRoot{p'},\allowbreak-e''\BareRoot{p''}\) etc. Hinc valor verus ipsius \(x\) erit
\SourceLeaf{039}
\[
=A-\alpha e\BareRoot{p}-\alpha'e'\BareRoot{p'}-\alpha''e''\BareRoot{p''}-\text{etc.}
\]
sive error valoris ipsius \(x\), in determinatione maxime idonea commissus, quem per \(Ex\) denotare convenit,
\[
=\alpha e\BareRoot{p}+\alpha'e'\BareRoot{p'}+\alpha''e''\BareRoot{p''}+\text{etc.}
\]
Perinde error valoris ipsius \(y\) in determinatione maxime idonea commissus, quem per \(Ey\) denotabimus, erit
\[
=\beta e\BareRoot{p}+\beta'e'\BareRoot{p'}+\beta''e''\BareRoot{p''}+\text{etc.}
\]
Valor medius quadrati \((Ex)^2\) invenitur
\[
=mmp(\alpha\alpha+\alpha'\alpha'+\alpha''\alpha''+\text{etc.})=mmp[\alpha\alpha]
\]
valor medius quadrati \((Ey)^2\) perinde \(=mmp[\beta\beta]\) etc., ut iam supra docuimus. Iam vero etiam valorem medium \textit{producti} \(Ex.Ey\) assignare licet, quippe qui invenitur
\[
=mmp(\alpha\beta+\alpha'\beta'+\alpha''\beta''+\text{etc.})=mmp[\alpha\beta]
\]
Concinne haec ita quoque exprimi possunt. Valores medii quadratorum \((Ex)^2,(Ey)^2\) etc. resp. aequales sunt productis ex \(\frac12mmp\) in quotientes differentialium partialium secundi ordinis
\[
\frac{dd\Omega}{d\xi^2},\quad\frac{dd\Omega}{d\eta^2}\quad\text{etc.}
\]
valorque medius producti talis, ut \(Ex.Ey\), aequalis est producto ex \(\frac12mmp\) in quotientem differentialem \(\frac{dd\Omega}{d\xi.d\eta}\): quatenus quidem \(\Omega\) tamquam functio indeterminatarum \(\xi,\eta,\zeta\) etc. consideratur.
\Article{29}
Designet \(t\) functionem datam linearem quantitatum \(x,y,z\) etc. puta sit
\[
t=fx+gy+hz+\text{etc.}+k
\]
Valor ipsius \(t\), e valoribus maxime plausibilibus ipsarum \(x,y,z\) etc. prodiens hinc erit \(=fA+gB+hC+\text{etc.}+k\), quem per \(K\) denotabimus. Qui si tamquam valor verus ipsius \(t\) adoptatur, error committitur, qui erit
